\begin{block}{Resumen}
  \begin{itemize}
    \item \textbf{Problema} La ecuación de Dirac quiral incorpora un ángulo
      en la dinámica del campo espinorial. \\
      El ángulo quiral induce un término de masa pseudoescalar. \\
      El lagrangiano de Yukawa se anula para las excitaciones levógiras,
      prediciendo neutrinos con masa nula.

    \item \textbf{Objetivo} Implementar un mecanismo seesaw quiral para generar
      dinámicamente un término de masa de Majorana.

    \item \textbf{Resultados} La estructura produce una excitación dextrógira
      estéril masiva y un campo activo levógiro con masa casi nula. \\
      El límite de ángulo quiral nulo recupera la teoría estándar.
  \end{itemize}
\end{block}
